% 第六章
\chapter{32 个点群与 230 个空间群的双值表示}

% 6.1 点群的双值表示
\section{点群的双值表示}

2.1 节考虑了三维转动群 $O(3)$ 的表示 $\mathcal{D}^{l}\{R(\alpha, \beta, \gamma)\}$，其中假定 $l$ 为整数。该表示的基是其分量为球谐函数 $Y^{m}_{l}(\theta, \phi)$（$-l \leqslant m \leqslant l$）的矢量。第二章列出的点群不可约表示（见表 2.2 与 2.3）可以认为是通过对这些表示的限制（subduction）得到的，即把表示 $\mathcal{D}^{l}\{R(\alpha, \beta, \gamma)\}$ 限制到某个点群的元素上（见定理 4.1.5 之下）。例如 Bethe (1929) 就是这样做的。点群的这些表示称为点群的\textit{单值表示}（single-valued representations），因为它们由三维转动群的、$l$ 为整数的表示 $\mathcal{D}^{l}\{R(\alpha, \beta, \gamma)\}$ 导出，而后者本身是单值的（即每个算符只有一个矩阵代表）。

我们所用的欧拉角 $\alpha$、$\beta$、$\gamma$ 的定义已在 2.1 节给出。对三维转动群的每个转动 $R(\alpha, \beta, \gamma)$，对每个整数值 $l$ 有且只有一个矩阵 $\mathcal{D}^{l}\{R(\alpha, \beta, \gamma)\}$；对任意给定的 $l$，矩阵 $\mathcal{D}^{l}\{R(\alpha, \beta, \gamma)\}$ 构成定义 1.3.4 意义下转动 $R(\alpha, \beta, \gamma)$ 群的一个表示。众所周知，球谐函数 $Y^{m}_{l}(\theta, \phi)$ 是量子力学角动量算符的本征函数：例如原子中电子波函数的角部分含有因子 $Y^{m}_{l}(\theta, \phi)$，且该电子轨道角动量的大小与 $z$ 分量分别为 $\sqrt{l(l+1)}\,\hbar$ 与 $m\hbar$。

然而，当在能带结构计算中引入相对论效应、或在晶体场理论中计入自旋—轨道耦合时，常常必须考虑电子自旋的存在。此时需要考虑的不是 $l$ 为整数的表示 $\mathcal{D}^{l}\{R(\alpha, \beta, \gamma)\}$，而是第二组表示 $\mathcal{D}^{j}\{R(\alpha, \beta, \gamma)\}$，其中 $j$ 为半奇整数。这种表示的基是旋量（spinor）而不是矢量。其中最简单的是 $\mathcal{D}^{\frac{1}{2}}\{R(\alpha, \beta, \gamma)\}$：
\begin{equation*}
\begin{array}{rl}
\mathcal{D}^{\frac{1}{2}}\{R(\alpha, \beta, \gamma)\} &= \pm\mathbf{u}\{R(\alpha, \beta, \gamma)\}\\
&= \pm\begin{pmatrix} e^{\frac{1}{2}i\alpha}\cos\frac{1}{2}\beta\, e^{\frac{1}{2}i\gamma} & e^{-\frac{1}{2}i\alpha}\sin\frac{1}{2}\beta\, e^{\frac{1}{2}i\gamma}\\ -e^{\frac{1}{2}i\alpha}\sin\frac{1}{2}\beta\, e^{-\frac{1}{2}i\gamma} & e^{\frac{1}{2}i\alpha}\cos\frac{1}{2}\beta\, e^{-\frac{1}{2}i\gamma} \end{pmatrix}.
\end{array}
\tag{6.1.1}
\end{equation*}
注意 $\mathcal{D}^{j}\{R(\alpha, \beta, \gamma)\}$ 与 $-\mathcal{D}^{j}\{R(\alpha, \beta, \gamma)\}$ 都必须包含在内，$\mathcal{D}^{j}\{R(\alpha, \beta, \gamma)\}$ 才能构成群。$\mathcal{D}^{j}\{R(\alpha, \beta, \gamma)\}$ 的矩阵元可写为
\begin{equation*}
\pm\mathcal{D}^{j}\{R(\alpha, \beta, \gamma)\}_{nm} = \exp(-\mathrm{i}m\alpha)\exp(-\mathrm{i}n\gamma)d^{j}(\beta)_{nm}
\tag{6.1.2}
\end{equation*}
其中 $d^{j}(\beta)_{nm}$ 由式 (2.1.6) 中把 $l$ 换成 $j$ 得到。第二章用过的因子 $C_{nm}$ 此处不需要，因为对这里用的旋量我们采用 Condon and Shortley (1935) 的相位因子，尽管对球谐函数我们用了不同于 Condon and Shortley 的相位因子。若让 $\alpha, \beta, \gamma$ 取遍所有允许值，使操作 $R(\alpha, \beta, \gamma)$ 包含三维纯转动群的所有成员，则矩阵 $\pm\mathbf{u}\{R(\alpha, \beta, \gamma)\}$ 将包含二维的全部幺正幺模矩阵，即群 $SU(2)$ 的全部成员。因此这些表示常被称为\textit{双值表示}（double-valued representations）。由转动群的这些双值表示，通过限制也可以得到点群的双值表示；点群的这种双值表示曾由 Bethe (1929) 与 Opechowski (1940) 考虑。

点群的任何成员或者是纯转动操作 $R(\alpha, \beta, \gamma)$，或者可以表示为某个这种纯转动操作与反演操作的乘积。于是一个点群可能完全由纯转动组成，也可能其一半的元素可以表示为反演操作与纯转动的乘积。后一种情形下，该点群或者同构于一个只含纯转动的点群，或者是一个只含纯转动的点群与含恒等元及反演的群 $\bar{1}\ (C_{i})$ 的直积。因此显然，实际上只需求出那些只含纯转动的点群的表示，其余点群的表示立刻随之得到。

设 $\mathbf{G}$ 是阶为 $|\mathbf{G}|$、由纯转动 $R(\alpha, \beta, \gamma)$ 组成的点群，则点群中的每个转动 $R(\alpha, \beta, \gamma)$ 对应于式 (6.1.1) 给出的两个矩阵 $\mathcal{D}^{\frac{1}{2}}\{R(\alpha, \beta, \gamma)\}$。这些矩阵构成一个 $2|\mathbf{G}|$ 阶的群，称为 $\mathbf{G}$ 的\textit{双群}（double group），记作 $\mathbf{G}^{\dagger}$。

\noindent\textbf{定义 6.1.1}\quad （Opechowski 1940）\ 阶为 $|\mathbf{G}|$ 的群 $\mathbf{G}$（它是三维转动群 $O(3)$ 的子群）的双群 $\mathbf{G}^{\dagger}$ 是 $2|\mathbf{G}|$ 阶的抽象群，其群乘法表与对应于 $\mathbf{G}$ 诸元素的 $SU(2)$ 的 $2|\mathbf{G}|$ 个矩阵的群乘法表相同。

虽然双群是用 $\mathcal{D}^{\frac{1}{2}}\{R(\alpha, \beta, \gamma)\}$ 定义的，用 $j$ 取任何其他半奇整数值的 $\mathcal{D}^{j}\{R(\alpha, \beta, \gamma)\}$ 也会得到完全相同的抽象群。用通常的方法、借助定理 1.3.8、1.3.10、1.3.11、1.3.13 与 1.3.15 可以求出这个群的不可约表示。$\mathbf{G}$ 是三维转动群的子群，而 $\mathbf{G}^{\dagger}$ 是 $SU(2)$（全体 $2\times2$ 幺正幺模矩阵群）的子群；记得 $SU(2)$ 与三维纯转动群之间存在二对一同态。

Opechowski (1940) 导出的一组规则使 $\mathbf{G}^{\dagger}$ 分解为类以及 $\mathbf{G}^{\dagger}$ 特征标表的计算得以简化。然而这些规则适用面有限：它们对求双群的特征标有用，但如果想超出这一点、导出实际的矩阵代表而不仅是特征标，就没有多大用处。在继续之前，我们只引用这些结果中较重要的几条，它们有时被称为“Opechowski 规则”。

\noindent\textbf{定理 6.1.1}\quad 点群 $\mathbf{G}$ 的不是 $\pi$ 转动类的每个类 $C_{i}$，恰对应于 $\mathbf{G}^{\dagger}$ 的两个类 $C'_{i}$ 与 $C''_{i}$。

\noindent\textbf{定理 6.1.2}\quad 点群 $\mathbf{G}$ 中是 $\pi$ 转动类的每个类 $C_{\pi}$，对应于 $\mathbf{G}^{\dagger}$ 的一个或两个类。若 $\mathbf{G}$ 中不存在绕垂直于 $C_{\pi}$ 中某转动之轴的轴的 $\pi$ 转动，则 $\mathbf{G}^{\dagger}$ 中有两个类对应于 $C_{\pi}$；若 $\mathbf{G}$ 中确实存在绕垂直于 $C_{\pi}$ 中某转动之轴的轴的 $\pi$ 转动，则 $\mathbf{G}^{\dagger}$ 中只有一个类对应于 $C_{\pi}$。

\noindent\textbf{定理 6.1.3}\quad $\mathbf{G}$ 的每个表示也是 $\mathbf{G}^{\dagger}$ 的表示，其中 $\chi(C''_{i}) = +\chi(C'_{i})$；它们称为 $\mathbf{G}$ 的单值表示。

\noindent\textbf{定理 6.1.4}\quad $\mathbf{G}^{\dagger}$ 的其余表示称为 $\mathbf{G}$ 的双值表示，满足 $\chi(C''_{i}) = -\chi(C'_{i})$。

习惯上把群 $\mathbf{G}^{\dagger}$ 的表示当作点群 $\mathbf{G}$ 本身的表示来考虑（如定理 6.1.3 与 6.1.4）；也就是说，若由元素 $R(\alpha, \beta, \gamma)$ 导出的 $\mathbf{G}^{\dagger}$ 两个元素的矩阵代表的特征标总是相同，就称它们为 $\mathbf{G}$ 的单值表示；若特征标总是相反，就称为 $\mathbf{G}$ 的双值表示。

表 6.1 给出各双群的特征标表与矩阵代表；表 6.2 给出双群 $432\ (O)$ 的群乘法表（(a) 无杠算符、(b) 带杠算符两种排法，识别带杠与无杠算符后恰与表 1.5、1.6 重合；元素上的横杠表示该矩阵为表 6.1 所示矩阵的 $-1$ 倍），表 6.3 给出双群 $422\ (D_{4})$ 的乘法表（只给出四分之一，其余由 $\overline{E}$ 与群中每个元素可交换得到），表 6.4 与 6.5 给出相关的小群识别。以下为原书各表扫描。

\begin{figure}[H]
\centering
\includegraphics[width=0.86\textwidth]{c6_61_p433.jpg}
\caption*{\textbf{表 6.1}\ 双值点群的特征标表与矩阵代表（原书第 420 页起扫描）}
\end{figure}
\begin{figure}[H]
\centering
\includegraphics[width=0.86\textwidth]{c6_61_p434.jpg}
\caption*{原书第 421 页}
\end{figure}
\begin{figure}[H]
\centering
\includegraphics[width=0.86\textwidth]{c6_61_p435.jpg}
\caption*{原书第 422 页}
\end{figure}
\begin{figure}[H]
\centering
\includegraphics[width=0.86\textwidth]{c6_61_p436.jpg}
\caption*{原书第 423 页}
\end{figure}
\begin{figure}[H]
\centering
\includegraphics[width=0.86\textwidth]{c6_61_p437.jpg}
\caption*{原书第 424 页}
\end{figure}
\begin{figure}[H]
\centering
\includegraphics[width=0.86\textwidth]{c6_61_p438.jpg}
\caption*{原书第 425 页}
\end{figure}
\begin{figure}[H]
\centering
\includegraphics[width=0.86\textwidth]{c6_61_p439.jpg}
\caption*{原书第 426 页}
\end{figure}
\begin{figure}[H]
\centering
\includegraphics[width=0.86\textwidth]{c6_61_p440.jpg}
\caption*{原书第 427 页}
\end{figure}
\begin{figure}[H]
\centering
\includegraphics[width=0.86\textwidth]{c6_61_p441.jpg}
\caption*{原书第 428 页}
\end{figure}
\begin{figure}[H]
\centering
\includegraphics[width=0.86\textwidth]{c6_61_p442.jpg}
\caption*{原书第 429 页}
\end{figure}
\begin{figure}[H]
\centering
\includegraphics[width=0.86\textwidth]{c6_61_p443.jpg}
\caption*{原书第 430 页}
\end{figure}
\begin{figure}[H]
\centering
\includegraphics[width=0.86\textwidth]{c6_61_p444.jpg}
\caption*{原书第 431 页}
\end{figure}
\begin{figure}[H]
\centering
\includegraphics[width=0.86\textwidth]{c6_61_p445.jpg}
\caption*{原书第 432 页}
\end{figure}
\begin{figure}[H]
\centering
\includegraphics[width=0.86\textwidth]{c6_61_p446.jpg}
\caption*{原书第 433 页}
\end{figure}
\begin{figure}[H]
\centering
\includegraphics[width=0.86\textwidth]{c6_61_p447.jpg}
\caption*{原书第 434 页}
\end{figure}
\begin{figure}[H]
\centering
\includegraphics[width=0.86\textwidth]{c6_61_p448.jpg}
\caption*{原书第 435 页}
\end{figure}
\begin{figure}[H]
\centering
\includegraphics[width=0.86\textwidth]{c6_61_p449.jpg}
\caption*{原书第 436 页}
\end{figure}
